\begin{textsample}{POS Dim 5 – global\_south\_2024\_09 – Score 16.00 – global\_south\_2024\_09/114428905.txt}  \label{ex:f5_pos_233}
UNITED NATIONS, Sep 19 : India \textbf{abstained} in the UN General Assembly on a \textbf{resolution} that demanded Israel end its unlawful presence in the Occupied Palestinian Territory within 12 months, with Delhi underlining that it is a strong advocate of"dialogue and diplomacy"and efforts should be made towards"building bridges","not furthering divides".
The 193-member General Assembly adopted the \textbf{resolution} on Wednesday, with 124 nations voting in \textbf{favour}, 14 against and 43 \textbf{abstentions}, including that by India.
The \textbf{resolution} was titled ' \textbf{Advisory} opinion of the \textbf{International} \textbf{Court} of Justice on the \textbf{legal} consequences arising from Israel ' s policies and practices in the Occupied Palestinian Territory, including East Jerusalem, and from the illegality of Israel ' s continued presence in the Occupied Palestinian Territory '.
Those \textbf{abstaining} included Australia, Canada, Germany, Italy, Nepal, Ukraine and the United Kingdom.
Nations voting against the \textbf{resolution} included Israel and the US.
Delivering the explanation of vote on the \textbf{resolution}, India ' s Permanent Representative to the UN Ambassador Parvathaneni Harish underlined India ' s abiding commitment to achieving @ @ @ @ @ @ @ @ @ @ issue and reiterated that only a two-state solution achieved through direct and meaningful negotiations between both sides will lead to enduring peace."India has \textbf{abstained} from today ' s vote.
We have been strong advocates of dialogue and diplomacy.
We believe that there is no other way to resolve conflicts,"he said.
Underlining that there are no winners in conflicts, Harish said the cost of conflict is human lives and destruction."A joint effort should be directed towards bringing the two sides closer, not drive them further apart.
We should strive towards building bridges, not furthering the divides,"he said.
India urged the General Assembly to make a"genuine effort"to strive for peace.
Harish emphasised India ' s steadfast commitment to \textbf{resolution} of the conflict and restoration of peace by bringing human suffering to an end."We will continue to be guided by this spirit.
To this end, we stand ready to continue our engagements towards achieving sustained peace,"he said.
Harish said the world @ @ @ @ @ @ @ @ @ @ than 11 months now that has led to the death of thousands including women and children and the largescale consequent humanitarian crisis.
Underlining that India ' s position on the conflict has been clear and consistent, Harish said Delhi \textbf{unequivocally} condemned the terror attacks on Israel on October 7 last year, and also condemned the loss of civilian lives in the conflict, called for an immediate ceasefire, and an immediate and \textbf{unconditional} release of all hostages.
He said India stands for unrestricted and sustained humanitarian assistance in the Gaza Strip.
The \textbf{resolution} adopted Wednesday demanded that"Israel brings to an end without delay its unlawful presence in the Occupied Palestinian Territory, which constitutes a wrongful \textbf{act} of a continuing character entailing its \textbf{international} responsibility, and do so no later than 12 months from the adoption of the present \textbf{resolution}."The Palestinian-drafted \textbf{resolution} also strongly \textbf{deplored} the continued and total disregard and breaches by the Government of Israel of its \textbf{obligations} under the Charter of the United Nations, \textbf{international} \textbf{law} and the relevant United Nations \textbf{resolutions}, and stressed @ @ @ @ @ @ @ @ @ @ security.
It \textbf{recognised} that Israel must be held to account for any \textbf{violations} of \textbf{international} \textbf{law} in the Occupied Palestinian Territory, including any \textbf{violations} of \textbf{international} humanitarian \textbf{law} and \textbf{international} human rights \textbf{law}, and that it"must bear the \textbf{legal} consequences of all its internationally wrongful \textbf{acts}, including by making reparation for the injury, including any damage, caused by such \textbf{acts}.
UN Secretary General Antonio Guterres, addressing reporters as the organisation prepares to host world leaders soon for the landmark Summit of the Future and the high-level General Assembly session said that the level of death and destruction in Gaza is by far the"largest that I ' ve ever seen in my mandate as Secretary-General"and the \textbf{violations} of \textbf{international} humanitarian \textbf{law} are"totally unacceptable.""I condemned as I did, and I feel that that condemnation was essential, the terror attacks by Hamas, but what ' s happening today in Gaza is totally unacceptable,"Guterres said.
The UN Office for the Coordination of Humanitarian Affairs ( OCHA ), citing @ @ @ @ @ @ @ @ @ @ October 7, 2023 and September 16 this year, least 41,226 Palestinians have been killed and 95,413 injured.
According to the Israeli military and official Israeli sources cited in the media, more than 1,542 Israelis and foreign nationals have been killed during this period, the majority on October 7 and its immediate aftermath. ( PTI )

% matched lemmas: abstain, abstention, act, advisory, court, deplore, favour, international, law, legal, obligation, recognise, resolution, unconditional, unequivocally, violation
\end{textsample}
